% 公共样式文件：只放宏包、页面、字体、标题和目录规则。
% 具体正文内容应放在 main.tex 或单独的内容块里。

\usepackage{ctex}
\usepackage{amsmath}
\usepackage{graphicx}
\usepackage{booktabs}
% 页面尺寸和页眉页脚控制
\usepackage{geometry}
\usepackage{fancyhdr}
% 章节标题与目录格式控制
\usepackage{titlesec}
\usepackage{setspace}
\usepackage{caption}
% 浮动体和超链接支持
\usepackage{float}
\usepackage{hyperref}
% 目录条目排版控制
\usepackage{titletoc}
% 空白页保持真正空白
\usepackage{emptypage}
% 带边框内容块
\usepackage[most]{tcolorbox}

% 按课程设计常见要求设置页边距
\geometry{left=3.18cm,right=3.18cm,top=2.54cm,bottom=2.54cm}
% 页眉高度留足，避免 fancyhdr 警告
\setlength{\headheight}{15pt}
% 正文统一使用 1.5 倍行距
\setstretch{1.5}

% 字体定义：中文宋体/黑体，西文 Times New Roman
\setCJKmainfont{SimSun}[AutoFakeBold=2.5]
\setCJKsansfont{SimHei}[AutoFakeBold=2.5]
\setmainfont{Times New Roman}
\setCJKfamilyfont{zhsong}{SimSun}[AutoFakeBold=2.5]
\setCJKfamilyfont{zhhei}{SimHei}[AutoFakeBold=2.5]
\renewcommand{\songti}{\CJKfamily{zhsong}}
\renewcommand{\heiti}{\CJKfamily{zhhei}}

% 一级标题：章，居中大号黑体
\titleformat{\chapter}[block]
  {\centering\zihao{-3}\bfseries\heiti}
  {\thechapter}
  {1em}
  {}
% 章标题和正文之间的间距
\titlespacing*{\chapter}{0pt}{-1.5\baselineskip}{0.5\baselineskip}

% 二级标题：节，四号黑体
\titleformat{\section}[block]
  {\zihao{4}\bfseries\heiti}
  {\thesection}
  {1em}
  {}

% 三级标题：小节，小四黑体
\titleformat{\subsection}[block]
  {\zihao{-4}\bfseries\heiti}
  {\thesubsection}
  {1em}
  {}

% 目录中章条目样式
\titlecontents{chapter}[0pt]
{\addvspace{6pt}\zihao{-4}\songti}
{\thecontentslabel\quad}
{}
{\titlerule*[5pt]{.}\thecontentspage}
% 目录中节条目样式
\titlecontents{section}[0pt]
{\addvspace{3pt}\zihao{-4}\songti}
{\thecontentslabel\quad}
{}
{\titlerule*[5pt]{.}\thecontentspage}
% 目录中小节条目样式
\titlecontents{subsection}[2em]
{\addvspace{2pt}\zihao{-4}\songti}
{\thecontentslabel\quad}
{}
{\titlerule*[5pt]{.}\thecontentspage}
% 目录深度控制到小节
\setcounter{tocdepth}{2}

% 图表标题格式
\DeclareCaptionFont{five}{\zihao{5}\bfseries\heiti}
\captionsetup{font=five,labelsep=space}